\section{Cantor space, gaps, and continuous images}\label{sec:cantor}
The Cantor set appears in the starting article. Its linked article
\emph{Compact space} gives the classical theorem that every nonempty compact
metric space is a continuous image of the Cantor set. We give the
arithmetical version for a dense rational metric domain with explicit finite
nets. An input point has a preimage computable from its name.

A binary stream $a$ is a fixed PR bit function with optional numerical
parameters. Write $a|n$ for its first $n$ bits, numbered from zero. A finite
word is its length together with its bit mask. The cylinder $[w]$ consists of
the streams with prefix $w$. Equality of streams means equality at each
index; apartness has a differing index as witness. Uniform continuity from
streams into a metric presentation means that agreement on the first
$L(k)$ bits guarantees output distance at most $2^{-k}$ for a supplied PR
function $L$.

\begin{theorem}[Clopen cylinders and finite prefix covers]\label{idx-cantor17}
Cylinders form a countable clopen basis. Distinct witnessed bits separate
points by complementary clopens. A connected subset contains at most one
point up to equality. There are no isolated points. Every finite list of
cylinders admits an exact cover test, with an explicit missed stream on
failure and a disjoint covering sublist on success.
\end{theorem}
\begin{proof}
Two cylinders intersect in the one with the longer prefix if their bits
are compatible, and otherwise are disjoint. The complement of $[w]$ is
the union of the other cylinders of length $|w|$. Finite intersections
and these complement formulas give the assertions about the product
topology whose basic restrictions fix finitely many individual bits:
any such restriction is a finite union of cylinders, taking the length
just past its largest restricted position. If two streams differ at $i$,
partition by bit $i$. Its two sides are finite unions of cylinders and
both meet the subset. Consequently, in a connected subset, for each fixed
$i$ the rationally decidable test of bit equality cannot fail, which proves
equality at every index. Flipping bit $n$ produces a distinct point in
$[a|n]$. Extending finite words by zeros also gives a countable dense list.

For words $W$ let $N$ be their maximum length, with $N=0$ for the empty
list. Enumerate all $2^N$ words of length $N$, asking whether each extends
an entry of $W$. A failing word followed by zeros misses every cylinder;
success proves coverage for any stream by reading exactly $N$ bits. Remove
duplicates and all words having a shorter listed prefix. This leaves the
same union, now disjoint. In the covering case each remaining $w$ accounts
for $2^{N-|w|}$ distinct length-$N$ words, so the retained words also
satisfy $\sum_w2^{-|w|}=1$. More generally, a PR predicate on words and a
given uniform hitting depth $N$ give a finite cover by scanning the words
of length at most $N$. Merely deciding the predicate does not supply this
depth.
\end{proof}

The $2^n$ zero-extended length-$n$ words give a finite net to precision
$n$ in the prefix topology. There is a complementary completeness
statement. Suppose a PR array of streams $a_i$ agrees in its first $k$ bits
whenever both row indices are at least the supplied $M(k)$. Define the
limit's bit $j$ to be bit $j$ of row $M(j+1)$. If $i\ge M(k)$ and $j<k$,
compare row $i$ with row $M(j+1)$ through the row indexed by
$\max(M(k),M(j+1))$. The first comparison preserves the first $k$ bits and
the second preserves the first $j+1$. Hence the limit has the required
convergence modulus. Uniqueness follows one bit at a time.

Similarly every represented continuous map from the interval to binary
streams is constant when for each output bit its evaluation on rationals
is PR and it is locally constant at every named input. Apply the cube's
bit-coloring theorem in Section~\ref{sec:metric-names} separately to each
bit. In particular the interval and Cantor set cannot be homeomorphic
by maps with these data, since the two constant streams have different
first bits.

\begin{theorem}[The middle-thirds embedding]\label{idx-cantor69}
Binary streams with the prefix topology and the named middle-thirds Cantor
set are uniformly homeomorphic. The latter is closed, perfect and nowhere
dense; it has finite interval covers of arbitrarily small total length.
The digit decoder works at the endpoints of all the deleted intervals.
\end{theorem}
\begin{proof}
For a word of length $n$ set
\[
 c(w)=\sum_{i<n}\frac{2w_i}{3^{i+1}},\qquad
 I_w=[c(w),c(w)+3^{-n}],\qquad C_n=\bigcup_{|w|=n}I_w .
\]
Children occupy the first and last thirds of their parent. Two different
level-$n$ intervals have distance at least $3^{-n}$: their first
different bit $i<n$ separates them by a middle gap of length
$3^{-(i+1)}$. For a rational $v$, distance to $C_n$ is the minimum of the
$2^n$ rational distances to the individual closed intervals; each distance
is $\max(0,c(w)-v,v-c(w)-3^{-n})$. It is 1-Lipschitz as a finite
minimum of 1-Lipschitz functions. For a regular name $x$ from
Section~\ref{sec:metric-names}, membership in the limiting set means
\[
          \forall n,q\quad \operatorname{dist}(x_q,C_n)\le 2^{-q}.
\]
A violation puts $x$ in the presented open complement: if the surplus is
$m>0$, a rational ball at $x_q$ of radius $2^{-q}+m/2$ contains $x$,
and is disjoint from $C_n$ by the distance estimate. Conversely enumerate
all rational balls with radius strictly below the center's distance to
some $C_n$. Membership in one yields a strict violation for a sufficiently
accurate sample. In particular membership in the Cantor set is the
negation of belonging to this open complement, since all individual
rational violation tests are decidable. This complement is the exterior
of $[0,1]$ together with the open middle thirds removed at the successive
stages. Indeed each finite-stage complement is exactly that finite list of
gaps together with the exterior; strict separation from its finite union
of closed intervals decides which gap by a sufficiently accurate sample.

Take $x_n=c(a|n)$. Subsequent sums lie in $I_{a|n}$ by the geometric sum,
so these left endpoints are a regular name with prefix error $3^{-n}$.
Each fixed interval $I_{a|n}$ contains the named point. Conversely,
given a name in the limiting set and a requested length $n$, take its
sample $x_h$ with $h=2n+3$ and select the word minimizing its distance
to $x_h$. For $n=0$ there is just the empty word. For $n>0$, its interval
is within $2^{-h}$ of the sample. At any $q\ge h$ the minimizing word
for $x_q$ must be the same: two different intervals would have distance
at most $2^{-h}+(2^{-h}+2^{-q})+2^{-q}\le4\cdot2^{-h}<3^{-n}$.
Thus at all such $q$ the selected interval is within $2^{-q}$ of
$x_q$. At earlier stages $p$, the triangle inequality bounds its distance
from $x_p$ by $2^{-p}+2\cdot2^{-q}$ for arbitrarily large $q$, and
error removal gives $2^{-p}$. This proves the interval contains the
named point. Uniqueness follows from the positive gap. Since a level
$n+1$ interval is contained in its level $n$ parent, the computed
words are compatible. Reading the last bit at each positive length
produces a single PR stream. The reconstructed point is within
$3^{-n}$ of its left endpoint at every level, hence equals the input.
The reverse composite preserves the unique prefix at every length.
Agreement of $n$ input bits bounds output displacement by $3^{-n}$;
input displacement in the real line strictly below $3^{-n}$ forces
the same $n$ output bits. These are moduli in the two directions,
and the estimates use no test of equality to a rational endpoint.

Closedness and the absence of isolated points give perfection, in the
sense that every Cantor point is a limit of other Cantor points: flip
bit $n$ to get distance between $3^{-(n+1)}$ and $3^{-n}$.
For any rational open interval take its midpoint $t$ and choose $n$
with $3^{-n}$ less than its radius. Decide whether $t$ belongs to a
level-$n$ interval. If not, a sufficiently small rational interval about
$t$ misses $C_n$ by a computed positive distance. If it does, the open
middle third of that level-$n$ interval lies inside the original interval
and misses $C_{n+1}$. Thus every basic open interval contains an open gap,
which proves nowhere density. Finally $C_n$ has an interval cover of total
length $(2/3)^n$. Choosing $n=2(k+1)$ gives total length strictly below
$2^{-k}$. If open covering intervals are required, enlarge each end by
$(2^{-k}-(2/3)^n)/2^{n+2}$. The enlarged total remains below $2^{-k}$.
\end{proof}

\begin{theorem}[A Cantor surjection from finite metric nets]\label{idx-cantor144}
Let $D$ be a nonempty PR rational metric domain. For each $k$ suppose
we have a nonempty finite PR list $c_{k,0},\ldots,c_{k,m_k-1}$ in $D$
such that every $u\in D$ is within $2^{-(k+4)}$ of a member.
Then binary streams map uniformly continuously onto the named completion
of $D$. Given any regular name $x$ one can compute a lifting stream.
\end{theorem}
\begin{proof}
Give proposals for stage $k$ a block of $b_k=\max(1,m_k)$ bits; this
simple choice ensures $2^{b_k}\ge m_k$. Read the block as an integer
modulo $m_k$ to propose a center $c_k$. Set $y_0=c_0$ and use the rule
\[
 y_{k+1}=\begin{cases}
 c_{k+1}&d(y_k,c_{k+1})\le2^{-(k+2)},\\
 y_k&\text{otherwise}.
 \end{cases}
\]
The test uses rational centers even when past proposals were rejected.
At any later stage the distance from $y_k$ is bounded by the finite
geometric sum of step bounds, hence by $2^{-(k+1)}$. Thus $(y_k)$
is a regular name. If two inputs agree through stage $k$'s block,
their outputs are within $2^{-k}$ by bridging through the common
$y_k$. The required prefix length is $L(k+1)=\sum_{j\le k}b_j$.

For a given target $x$, choose the first center within $2^{-(k+4)}$ of
the rational sample $x_{k+4}$ by scanning the supplied finite list.
It is within $2^{-(k+3)}$ of the named point. Two consecutive chosen
centers have distance at most $2^{-(k+3)}+2^{-(k+4)}<2^{-(k+2)}$.
Thus every proposal in this stream is accepted, and the output equals
$x$ by the error estimate. To compute bit $i$ of the stream, find the
unique $k$ with $L(k)\le i<L(k+1)$, where $L(0)=0$. Scanning at most
$i+1$ blocks suffices since every block has positive length. Output
the requested bit in the padded representation of the chosen center
index. This proves the lifting claim without a choice principle.
Different names for the same point may give different lifting streams.
\end{proof}

For instance the cube grids of Section~\ref{sec:metric-names} give these
nets uniformly in the dimension. The Cantor set therefore maps onto every
such cube, including the whole interval. This construction permits redundant
representations of a target, as its surjectivity proof requires. Interleaving
bits is itself a homeomorphism from a pair of streams to one stream: undo
it by reading even and odd positions. Finally no fixed PR double array lists
all PR streams, since $n\mapsto1-a_{n,n}$ differs from row $n$ at position
$n$. This is a diagonal theorem scheme about uniform enumerations.

\paragraph{Lean implementation.}
Represent $w$ by a natural bit mask and its length, with prefix, append,
test-bit and truncation. The coverage certificate is a bounded
\texttt{all} over masks below $2^N$, with a bounded \texttt{any} over the
input list. Define $3^n c(w)$ by appending digits to a base-three integer.
Prove the interval gap by the first differing bit, computed by bounded
search. Derive the decoder's universal containment only after proving the
rational nearest-interval estimate for two fixed precisions. For the
surjection define the offset $L$ by primitive recursion and store the
current rational center as the sole state of the acceptance recursion.
The quantitative invariant at two fixed stages is the rational geometric
sum bound. In the lifting proof the invariant that the accepted center
equals the proposed center at the current stage is decidable. The theorem
schemes have only numerical inputs once the stream, net and target-stage
recipes are fixed.
